% Generated draft: requires mathematical and source review.
\documentclass[11pt]{article}
\usepackage[margin=25mm]{geometry}
\usepackage{fontspec}
\setmainfont{Latin Modern Roman}
\usepackage{amsmath,amssymb,mathtools,unicode-math}
\setmathfont{Latin Modern Math}
\usepackage{enumitem,needspace,xurl,hyperref}
\setlength{\parindent}{0pt}
\setlength{\parskip}{5pt}
\setlength{\emergencystretch}{3em}
\newcommand{\N}{\mathbb{N}}
\newcommand{\Z}{\mathbb{Z}}
\newcommand{\Q}{\mathbb{Q}}
\newcommand{\R}{\mathbb{R}}
\newcommand{\C}{\mathbb{C}}
\newcommand{\F}{\mathbb{F}}
\newcommand{\A}{\mathbb{A}}
\newcommand{\PP}{\mathbb{P}}
\newcommand{\E}{\mathbb{E}}
\newcommand{\eps}{\varepsilon}
\newcommand{\dd}{\,\mathrm{d}}
\begin{document}
\section*{UnsolvedMath ID: 30004503}
\textbf{Source problem code:} catalogue entry
\subsection{Definitions and mathematical statement}
The supplied record concerns irreducible supercuspidal representations of \(G=\operatorname{GL}_n(F)\). It does not specify the local field \(F\), the admissible values of \(n\), the representation category, or whether representations are required to be unitary. Write \(X_n(F)\) for the intended set of isomorphism classes, once these choices are fixed.

A pseudometric on a set \(X\) is a function \(d:X\times X\to[0,\infty)\) such that, for all \(x,y,z\in X\),
\[
d(x,x)=0,\qquad d(x,y)=d(y,x),\qquad d(x,z)\le d(x,y)+d(y,z).
\]
Distinct elements may have distance zero. Let \(d_{\mathrm{RS}}\) denote the local Rankin--Selberg pseudometric intended by the source. Its defining formula and normalization are absent from the supplied record; the axioms above do not identify that pseudometric.

For a center \(\pi\in X_n(F)\) and a real radius \(r\ge0\), the two possible ball conventions are
\[
B_r^{<}(\pi)=\{\sigma\in X_n(F):d_{\mathrm{RS}}(\pi,\sigma)<r\},
\qquad
B_r^{\le}(\pi)=\{\sigma\in X_n(F):d_{\mathrm{RS}}(\pi,\sigma)\le r\}.
\]
The record does not choose between them or specify the intended range of radii. In particular, the distinction matters at \(r=0\).

\textbf{Question, conditional on the missing specifications.} For the intended local field, rank, representation space, Rankin--Selberg pseudometric, and ball convention, determine
\[
\mu_{\mathrm{Pl}}\bigl(B_r(\pi)\bigr)
\]
as a function of the center \(\pi\) and the allowed radius \(r\), where \(\mu_{\mathrm{Pl}}\) is the intended Plancherel measure. For the usual group Plancherel interpretation, one must specify the relevant unitary representation space and the normalization of Haar measure on \(G\). The record does not explain how its supercuspidal representation space is related to that measure, nor does it supply a measurability assertion for these balls. Consequently, this is a formulation schema preserving the supplied question, rather than a fully specified measure-computation problem.
\subsection{Short English statement}
The question asks for the Plancherel measure of a ball centered at a supercuspidal representation of \(\operatorname{GL}_n(F)\), using the local Rankin--Selberg pseudometric. The supplied record leaves the pseudometric formula, representation domain, measure normalization, and ball boundary convention unspecified. These choices require clarification before a precise numerical or functional answer is determined.
\subsection{Sources}
\begin{itemize}
\item Local UnsolvedMath registry: supplied complete source record, permanent ID \(30004503\), source problem code \texttt{catalogue entry}, title ``Plancherel Volumes of Rankin--Selberg Metric Balls.'' Project registry references are \texttt{data/unsolvedmath/problems.json} and \texttt{catalog/problems.json}; the supplied material does not identify which file contains this entry.
\item Supplied report citation: \emph{Automorphic Forms and Arithmetic} (2021), Oberwolfach Reports, pp.~1317--1383, DOI \texttt{10.4171/owr/2020/26}. \url{https://doi.org/10.4171/owr/2020/26}. The background describes Oberwolfach Reports as records of mathematical workshops and their open-problem discussions. No individual proposer is supplied; the proposed year is recorded as 2020.
\item Supplied literature source: ``Oberwolfach Report 26/2020,'' \url{https://doi.org/10.4171/owr/2020/26}. The literature assessment checked on \texttt{2026-08-22} labels the problem open and describes the Plancherel-measure question as a natural question with no cited computation resolving it. This is supplied dated curation derived from the OWR material, not an independently verified literature assessment; its time-sensitive status and cited claims require independent verification.
\item The registry assigns the statement status \texttt{corrected\_verified} and says that mathematical content was preserved while notation, typography, omitted definitions, or duplicate clauses were repaired. This is a supplied editorial claim, not mathematical verification of this draft.
\end{itemize}
\end{document}
